\section{Rumah Susun}

\subsection{Definisi dan Landasan Hukum}
Rumah Susun adalah bangunan gedung bertingkat yang dibangun dalam suatu lingkungan yang terbagi dalam bagian-bagian yang distrukturkan secara fungsional, baik dalam arah horizontal maupun vertikal yang mana merupakan satuan-satuan yang masing-masing dapat dimiliki dan digunakan secara terpisah. Definisi dan regulasi utama mengenai rumah susun diatur dalam Undang-Undang Republik Indonesia Nomor 20 Tahun 2011 tentang Rumah Susun. Undang-undang ini hadir sebagai respon atas meningkatnya kebutuhan hunian di daerah perkotaan yang memiliki keterbatasan lahan (Sutedi, 2010).

Tujuan pembangunan rumah susun tidak hanya sekadar menyediakan tempat tinggal, tetapi juga untuk meningkatkan efisiensi penggunaan tanah, tata ruang, dan menciptakan lingkungan pemukiman yang lengkap, serasi, dan seimbang. Menurut Budiharjo dan Sujarto (2013) dalam bukunya \textit{Kota Berkelanjutan}, konsep hunian vertikal merupakan solusi paling rasional untuk mengatasi \textit{urban sprawl} (perembetan kota) yang tidak terkendali akibat tingginya arus urbanisasi.

\subsection{Klasifikasi Rumah Susun}
Berdasarkan peruntukan dan pengelolaannya, rumah susun dapat diklasifikasikan menjadi beberapa kategori. Menurut Undang-undang Nomor 20 Tahun 2011 tentang Rumah Susun\footnote{Undang-undang Nomor 20 Tahun 2011 tentang Rumah Susun (Lembaran Negara Republik Indonesia Tahun 2011 Nomor 108)}, klasifikasi tersebut meliputi berikut.

\begin{enumerate}
    \item \textbf{Rumah Susun Umum} \\
    Rumah susun yang diselenggarakan untuk memenuhi kebutuhan rumah bagi Masyarakat Berpenghasilan Rendah (MBR).
    \item \textbf{Rumah Susun Khusus} \\
    Rumah susun yang diselenggarakan untuk memenuhi kebutuhan khusus, seperti asrama mahasiswa atau hunian bagi lansia.
    \item \textbf{Rumah Susun Negara} \\
    Rumah susun yang dimiliki negara dan berfungsi sebagai tempat tinggal atau hunian sarana pembinaan keluarga serta penunjang pelaksanaan tugas pejabat dan/atau pegawai negeri.
    \item \textbf{Rumah Susun Komersial} \\
    Rumah susun yang diselenggarakan untuk mendapatkan keuntungan, yang sering dikenal dengan istilah {Apartemen} atau Kondominium. Objek penelitian ini, yaitu The Jarrdin Cihampelas yang mana termasuk dalam kategori Rusunami (Rumah Susun Sederhana Milik) yang dalam perkembangannya sering kali beroperasi dengan mekanisme pasar komersial.
\end{enumerate}

\subsection{Struktur Kepemilikan}
Salah satu karakteristik unik dari rumah susun yang membedakannya dengan rumah tapak (\textit{landed house}) adalah konsep kepemilikannya. Pemilik unit rumah susun tidak hanya memiliki unit huniannya secara individual, tetapi juga memiliki hak atas bagian lain dari gedung tersebut secara kolektif.

Menurut Harsono (1999) dalam bukunya berjudul \textit{Hukum Agraria Indonesia, Himpunan Peraturan-Peraturan Hukum Tanah}, kepemilikan dalam rumah susun meliputi tiga elemen tak terpisahkan sebagai berikut.
\begin{enumerate}
    \item {Hak Milik atas Satuan Rumah Susun (HMSRS)} yaitu
    hak kepemilikan yang bersifat perseorangan dan terpisah (ruang unit hunian).
    \item {Bagian bersama} yaitu 
    Bagian rumah susun yang dimiliki secara tidak terpisah untuk pemakaian bersama dalam kesatuan fungsi gedung (contoh: pondasi, kolom, dinding pemikul, atap, dan tangga).
    \item {Benda bersama} yaitu 
    Benda yang bukan merupakan bagian rumah susun melainkan bagian yang dimiliki bersama untuk pemakaian bersama (contohnya, instalasi listrik, tempat parkir, taman, dan bangunan pertokoan).
    \item {Tanah bersama} yaitu
    Sebidang tanah yang digunakan atas dasar hak bersama secara tidak terpisah yang di atasnya berdiri rumah susun.
\end{enumerate}



% \subsection{Dinamika Sosial dan Potensi Kerawanan}
% Tinggal di rumah susun membawa perubahan pola interaksi sosial dari horizontal menjadi vertikal. Marlina (2008) dalam \textit{Panduan Perancangan Bangunan Komersial} menyebutkan bahwa karakteristik penghuni apartemen atau rumah susun cenderung lebih individualis dan memiliki privasi yang tinggi dibandingkan penghuni perumahan tapak.

% Kondisi ini memiliki dua sisi mata uang. Di satu sisi, privasi penghuni sangat terjaga. Namun, di sisi lain, tingginya tingkat anonimitas (ketidaktahuan antar tetangga) dan aksesibilitas yang sering kali longgar pada apartemen sewa harian menciptakan celah keamanan. Kurangnya kontrol sosial (\textit{social control}) dari sesama penghuni memudahkan terjadinya penyalahgunaan unit untuk aktivitas ilegal, seperti prostitusi terselubung atau peredaran narkoba, karena pelaku merasa terlindungi oleh dinding privasi dan kurangnya interaksi tetangga (Soekanto, 2012). Hal inilah yang mendasari perlunya sistem pengawasan berbasis teknologi, seperti \textit{data capture} dan deteksi anomali, untuk membantu manajemen gedung dalam memitigasi risiko keamanan yang muncul akibat karakteristik sosial hunian vertikal.